\lecture{20}{May 12.}{}
\begin{theorem}[Arzela-Ascoli's theorem: sequential form]
    Let \((K, d)\) be a compact metric space. If \(\left\{ f_n \right\}_{n=1}^{\infty} \subseteq C(K) \) is uniformly bounded and equi-continuous, then \(\left\{ f_n \right\} \) has a subsequence which converges uniformly on \(K\) to a function in \(C(K)\).   
\end{theorem}

\begin{proof}
\hfill
    \begin{itemize}
        \item \textbf{Step 1. Compact metric spaces are separable.} For each \(k \in \mathbb{N}\), the collection of balls 
        \[
            \left\{ B\left( x, \frac{1}{k} \right): x \in K  \right\} 
        \]
        is an open cover of \(K\). By compactness, choose finitely many points \(x_{k, 1}, x_{k, 2}, \dots , x_{k, N_k} \in K\) s.t. 
        \[
            K = \bigcup_{j=1}^{N_k} B \left( x_{k, j}, \frac{1}{k} \right).  
        \]  
        Set 
        \[
            S_k \coloneqq \left\{ x_{k, 1}, x_{k, 2}, \dots , x_{k, N_k} \right\}, \quad S \coloneqq \bigcup_{k=1}^{\infty} S_k.  
        \]
        Then \(S\) is countable. We claim that \(S\) is dense in \(K\). Indeed, given \(x \in K\) and \(\varepsilon > 0\), choose \(k \in \mathbb{N}\) so large that \(\frac{1}{k} < \varepsilon\). Since \(K = \bigcup_{j=1}^{N_k} B(x_{k, j}, \frac{1}{k}) \), there exists \(j\) s.t. \(x \in B(x_{k, j}, \frac{1}{k})\). Hence, 
        \[
            d(x, x_{k, j}) < \frac{1}{k} < \varepsilon,
        \]          
        and \(x_{k, j} \in S\). Therefore, \(S\) is dense in \(K\). 
        \item \textbf{Step 2. A diagonal subsequence converges on the dense set.} List the countable dense set as \(S = \left\{ p_1, p_2, \dots  \right\} \). If \(S\) is finite, we repeat its elements so that it is still written as an infinite sequence. 
        
        Since \(\left\{ f_n \right\} \) is uniformly bounded, there exists a constant \(M > 0\) s.t. 
        \[
            \vert f_n(x) \vert \le M \quad \forall n \in \mathbb{N} \text{ and } \forall x \in K.  
        \]  
        In particular, for each fixed \(p_j \in S\), the numerical sequence \(\left\{ f_n(p_j) \right\}_{n=1}^{\infty}  \) is bounded in \(\mathbb{R}\). We now extract subsequences one point at a time. 

        First consider the point \(p_1\). Since \(\left\{ f_n(p_1) \right\} \) is bounded in \(\mathbb{R} \), the Bolzano-Weierstrass theorem gives a subsequence \(\left\{ f_n^{(1)} \right\}_{n=1}^{\infty}  \) of \(\left\{ f_n \right\} \) s.t. \(\left\{ f_n^{(1)} (p_1) \right\}_{n=1}^{\infty}  \) converges in \(\mathbb{R}\). 

        Next, consider \(p_2\). The sequence \(\left\{ f_n^{(1)}(p_2) \right\} \) is bounded again, so it has a convergent subsequence. Thus, we may choose a subsequence \(\left\{ f_n^{(2)} \right\}_{n=1}^{\infty}  \) of \(\left\{ f_n^{(1)} \right\}_{n=1}^{\infty}  \) s.t. \(\left\{ f_n^{(2)}(p_2) \right\}_{n=1}^{\infty}  \) converges. Since \(\left\{ f_n^{(2)} \right\} \) is a subsequence of \(\left\{ f_n^{(1)} \right\} \), the sequence \(\left\{ f_n^{(2)}(p_1) \right\} \) also converges. 

        Continuing inductively, suppose that for some \(r \ge 1\) we have constructed a subsequence \(\left\{ f_n^{(r)} \right\}_{n=1}^{\infty}  \) s.t. \(\left\{ f_n^{(r)} (p_j) \right\}_{n=1}^{\infty}  \) converges for each \(j = 1, \dots ., r\). Since the sequence \(\left\{ f_n^{(r)} (p_{r+1}) \right\}_{n=1}^{\infty}  \) is bounded in \(\mathbb{R}\), it has a convergent subsequence. Choose such a subsequence and call it \(\left\{ f_n^{(r+1)} \right\}_{n=1}^{\infty} \). Then \(\left\{ f_n^{(r+1)} (p_{n+1}) \right\}_{n=1}^{\infty} \) converges. Moreover, because \(\left\{ f_n^{(r+1)} \right\} \) is a subsequence of \(\left\{ f_n^{(r)} \right\} \), convergence at the earlier point \(p_1, \dots , p_r\) is preserved. 

        In this way we obtain a nested seuqence of sequences:
        \[
            \left\{ f_n^{(1)} \right\} \supseteq \left\{ f_n^{(2)} \right\} \supseteq \left\{ f_n^{(3)} \right\} \supseteq \dots ,   
        \]
        where the \(r\)-th subsequence \(\left\{ f_n^{(r)} \right\} \) converges at the first \(r\) points \(p_1, \dots , p_r\). We now define the diagonal subsequence by 
        \[
            g_n \coloneqq f_n^{(n)}, \quad n = 1, 2, 3, \dots. 
        \]    
        This is a subsequence of the original sequence \(\left\{ f_n \right\} \). We claim that for every fixed \(p_j \in S\), the numerical sequence \(\left\{ g_n (p_j) \right\}_{n=1}^{\infty}  \) converges.  

        Indeed, fix \(j \in \mathbb{N}\). For every \(n \ge j\), the term \(g_n = f_n^{(n)}\) belongs to the subsequence \(\left\{ f_m^{(j)} \right\}_{m=1}^{\infty} \), because the subsequences are nested. Hence the tail \(\left\{ g_n (p_j) \right\}_{n \ge j} \) is a subsequence of the convergent sequence \(\left\{ f_m^{(j)} (p_j) \right\}_{m=1}^{\infty}  \). Therefore, \(\left\{ g_n (p_j) \right\}_{n \ge j} \) converges. Adding or removing finitely many terms does not affect convergence, so \(\left\{ g_n (p_j) \right\}_{n=1}^{\infty}  \) converges. Since \(j\) was arbitrary, the diagonal subsequence \(\left\{ g_n \right\} \) converges pointwise on the dense set \(S\). 

        \item \textbf{Step 3. Equicontinuity upgrades pointwise convergence on \(S\) to uniform Cauchy convergence on \(K\).} Let \(\varepsilon > 0\). By equicontinuity, choose \(\delta > 0\) such that 
        \[
            d(x, y) < \delta \implies \vert g_n(x) - g_n(y) \vert < \frac{\varepsilon}{3} 
        \]
        for every \(n \in \mathbb{N}\). Choose \(k \in \mathbb{N}\) so large that \(\frac{1}{k} < \delta\). Since \(S_k\) is finite and \(\left\{ g_n(s) \right\} \) converges for each \(s \in S_k\), there exists \(N \in \mathbb{N}\) s.t. 
        \[
            m, n \ge N \implies \vert g_n(s) - g_m(s) \vert < \frac{\varepsilon}{3} \quad \text{for every } s \in S_k.  
        \]       
        Fix any \(x \in K\), choose \(s \in S_k\) with \(d(x, s) < \frac{1}{k} < \delta\). Then, for \(m, n \ge N\),
        \begin{align*}
            \vert g_n(x) - g_m(x) \vert &\le \vert g_n(x) - g_n(s)  \vert + \vert g_n(s) - g_m(s) \vert + \vert g_m(s) - g_m(x) \vert \\
            &< \frac{\varepsilon}{3} + \frac{\varepsilon}{3} + \frac{\varepsilon}{3} = \varepsilon.    
        \end{align*}    
        Hence, \(\left\{ g_n \right\} \) is Cauchy in \(\left( C(K), \lVert \cdot \rVert_{\infty }  \right) \). Since \(C(K)\) is complete, there exists \(g \in C(K)\) such that 
        \[
            \lVert g_n - g \rVert_{\infty} \to 0.  
        \]    
        Thus, \(g_n \to g\) uniformly on \(K\), as desired.  
    \end{itemize}
\end{proof}

\begin{corollary}[Compactness criterion in \(C(K)\)]
    A subset \(\mathcal{F} \subseteq C(K)\) is compact if and only if it is closed, pointwise bounded, and equicontinuous. Equivalently, since pointwise boundedness plus equicontinuity implies uniform boundedness, compact subsets of \(C(K)\) are precisely the closed, uniformly bounded, equicontinuous subsets.  
\end{corollary}

\begin{proof}
    First suppose \(\mathcal{F}\) is compact. Since \(C(K)\) is a metric space, \(\mathcal{F} \) is closed. The evaluation map 
    \[
        E_x : C(K) \to \mathbb{R}, \quad E_x(f) \coloneqq f(x),
    \]   
    is continuous for each \(x \in K\) because
    \[
        \vert E_x(f) - E_x(g) \vert = \vert f(x) - g(x) \vert \le \sup_{t \in K} \vert f(t) - g(t) \vert = \lVert f - g \rVert_{\infty}. 
    \]
    Hence, \(E_x(\mathcal{F})\) is compact in \(\mathbb{R}\), and therefore bounded. Thus, \(\mathcal{F} \) is pointwise bounded.

    It remains to prove equicontinuity. Let \(\varepsilon > 0\). Since \(\mathcal{F} \) is compact in the metric induced by \(\lVert \cdot \rVert_\infty  \), it is totally bounded. Thus, there exist \(f_1, \dots , f_L \in \mathcal{F}\) such that for every \(f \in \mathcal{F}\), 
    \[
        \lVert f - f_i \rVert_\infty < \frac{\varepsilon}{3} 
    \]         
    for some \(i\). Since each \(f_i \in \mathcal{F} \subseteq C(K)\) is uniformly continuous and there are only finitely many of them, there exists \(\delta > 0\) s.t. 
    \[
        d(x, y) < \delta \implies \vert f_i(x) - f_i(y) \vert < \frac{\varepsilon}{3} 
    \]   
    for every \(i = 1, 2, \dots , L\). If \(f \in \mathcal{F}\) and \(i\) is chosen with \(\lVert f - f_i \rVert_\infty < \frac{\varepsilon}{3} \), then \(d(x, y) < \delta\) implies  
    \[
        \vert f(x) - f(y) \vert \le \vert f(x) - f_i(x) \vert + \vert f_i(x) - f_i(y) \vert + \vert f_i(y) - f(y) \vert < \varepsilon.    
    \]    
    Thus \(\mathcal{F}\) is equicontinuous. 

    Conversely, suppose \(\mathcal{F} \) is closed, pointwise bounded, and equicontinuous. Let \(\left\{ f_n \right\} \) be any sequence in \(\mathcal{F}\). By the preceding lemma, \(\mathcal{F} \) is uniformly bounded. The sequential Arzela-Ascoli theorem therefore gives a subsequence \(\left\{ f_{n_j} \right\} \) converging uniformly to some \(f \in C(K)\). Since \(\mathcal{F}\) is closed, \(f \in \mathcal{F}\). Hence every sequence in \(\mathcal{F} \) has a convergent subsequence with limit in \(\mathcal{F}\). Because \(C(K)\) is a metric space, sequential compactness is equivalent to compactness. Thus, \(\mathcal{F} \) is compact.            
\end{proof}

\begin{corollary}
    Let \(\left\{ f_n \right\} \subseteq C(K) \) be pointwise bounded and equicontinuous. Then \(\left\{ f_n \right\} \) has a subsequence which converges uniformly on \(K\). In particular, the same conclusion holds if the sequence is uniformly bounded and equicontinuous.
\end{corollary}
\begin{proof}
    By \autoref{lm: Pointwise boundedness becomes uniform boundedness} and Arzela-Ascoli Theorem, this is true. 
\end{proof}

\begin{eg}
    Let \(K\) be a finite set with \(n\) points, equipped with the discrete metric 
    \[
        d(x, y) = \begin{dcases}
            0, &\text{ if } x = y ;\\
            1, &\text{ if } x \neq y.
        \end{dcases}
    \]  
    Show that \(C(K)\) may be identified with \(\mathbb{R} ^n\), and that the supremum norm on \(C(K)\) is equivalent to the usual Euclidean norm on \(\mathbb{R} ^n\). Conclude that the Bolzano-Weierstrass theorem is a special case of Arzela-Ascoli theorem.     
\end{eg}

\begin{explanation}
    Let 
\[
K = \left\{ x_1, x_2, \dots , x_n \right\}
\]
be a finite set equipped with the discrete metric
\[
d(x,y)=
\begin{dcases}
0,&x=y,\\
1,&x\neq y.
\end{dcases}
\]

Since \(K\) is finite and discrete, every subset of \(K\) is open. Hence every function \(f:K\to\mathbb R\) is continuous, so
\[
C(K)=\{f:K\to\mathbb R\}.
\]

Define
\[
T:C(K)\to\mathbb R^n
\]
by
\[
T(f)=\bigl(f(x_1),f(x_2),\dots,f(x_n)\bigr).
\]
The map \(T\) is linear and bijective, since a function on \(K\) is completely determined by its values at the points \(x_1,\dots,x_n\). Thus \(C(K)\) may be identified with \(\mathbb R^n\).

Now let \(f\in C(K)\). The supremum norm on \(C(K)\) is
\[
\|f\|_\infty
=
\sup_{x\in K}|f(x)|.
\]
Because \(K\) has finitely many points,
\[
\|f\|_\infty
=
\max_{1\le i\le n}|f(x_i)|.
\]
Under the identification with \(\mathbb R^n\), this is exactly the max norm
\[
\|(a_1,\dots,a_n)\|_{\max}
=
\max_{1\le i\le n}|a_i|.
\]

We now compare this norm with the Euclidean norm
\[
\|(a_1,\dots,a_n)\|_2
=
\left(\sum_{i=1}^n a_i^2\right)^{1/2}.
\]
For every \(a=(a_1,\dots,a_n)\in\mathbb R^n\),
\[
\|a\|_{\max}\le \|a\|_2
\]
since
\[
|a_i|^2\le \sum_{j=1}^n a_j^2
\]
for each \(i\). Also,
\[
\|a\|_2^2
=
\sum_{i=1}^n a_i^2
\le
\sum_{i=1}^n \|a\|_{\max}^2
=
n\|a\|_{\max}^2,
\]
so
\[
\|a\|_2\le \sqrt n\,\|a\|_{\max}.
\]
Hence the supremum norm and the Euclidean norm are equivalent on \(\mathbb R^n\).

Finally, observe that every family of functions on \(K\) is equicontinuous. Indeed, if \(\varepsilon>0\), choose \(\delta=\frac12\). Then
\[
d(x,y)<\delta
\implies x=y,
\]
and therefore
\[
|f(x)-f(y)|=0<\varepsilon
\]
for every function \(f\in C(K)\).

Thus, in this setting, the Arzelà--Ascoli theorem states that every uniformly bounded sequence in \(C(K)\cong\mathbb R^n\) has a uniformly convergent subsequence. Since uniform convergence corresponds to convergence in the max norm, and this norm is equivalent to the Euclidean norm, we conclude that every bounded sequence in \(\mathbb R^n\) has a convergent subsequence.

This is precisely the Bolzano--Weierstrass theorem.
\end{explanation}

\section{Application: Compactness of Integral Operators}
We now give an important application of the Arzela-Ascoli theorem. Let \(K : [a, b] \times [a, b] \to \mathbb{R}\) be continuous. For each \(f \in C([a, b])\), define 
\[
    (Tf)(x) \coloneqq \int _a^b K(x, y) f(y) \, \mathrm{d} y, \quad x \in [a, b]. 
\]  
We shall first check that \(Tf\) is continuous, so that \(T\) is indeed a map from \(C([a, b])\) into itself. Then we shall prove that \(T\) is compact with respect to the supremum norm. 

\begin{proposition}
Let \(K : [a,b]\times[a,b]\to\mathbb{R}\) be continuous, and define
\[
    (Tf)(x)\coloneqq \int_a^b K(x,y)f(y)\,dy,
    \qquad x\in[a,b].
\]
Then for every \(f\in C([a,b])\), the function \(Tf\) is continuous on \([a,b]\).
\end{proposition}

\begin{proof}
Fix \(f\in C([a,b])\). Since \(f\) is continuous on the compact interval \([a,b]\), it is bounded. Hence there exists \(M>0\) such that
\[
    |f(y)|\le M
    \qquad \forall y\in[a,b].
\]

Let \(\{x_n\}_{n=1}^\infty \subseteq [a,b]\) satisfy \(x_n\to x\in[a,b]\). We shall show that
\[
    (Tf)(x_n)\to (Tf)(x).
\]

For each \(n\),
\[
\begin{aligned}
    |(Tf)(x_n)-(Tf)(x)|
    &= \left|
        \int_a^b K(x_n,y)f(y)\,dy
        -
        \int_a^b K(x,y)f(y)\,dy
    \right| \\
    &= \left|
        \int_a^b \bigl(K(x_n,y)-K(x,y)\bigr)f(y)\,dy
    \right| \\
    &\le
    \int_a^b
        |K(x_n,y)-K(x,y)|\,|f(y)|
    \,dy \\
    &\le
    M\int_a^b |K(x_n,y)-K(x,y)|\,dy.
\end{aligned}
\]

Since \(K\) is continuous on the compact set
\(
[a,b]\times[a,b],
\)
it is uniformly continuous. Therefore,
\[
    \sup_{y\in[a,b]}
    |K(x_n,y)-K(x,y)|
    \to 0
    \qquad \text{as } n\to\infty.
\]
Consequently,
\[
\begin{aligned}
    \int_a^b |K(x_n,y)-K(x,y)|\,dy
    &\le
    (b-a)
    \sup_{y\in[a,b]}
    |K(x_n,y)-K(x,y)| \\
    &\to 0.
\end{aligned}
\]

It follows that
\[
    |(Tf)(x_n)-(Tf)(x)|\to 0,
\]
and hence
\[
    (Tf)(x_n)\to (Tf)(x).
\]

Therefore \(Tf\) is continuous on \([a,b]\), i.e.
\[
    Tf\in C([a,b]).
\]
\end{proof}

\begin{theorem}[Compactness of integral operators]
    Let \(K : [a, b] \times [a, b] \to \mathbb{R} \) be continuous. Define \(T: C([a, b]) \to C([a, b])\) by 
    \[
        (Tf)(x) \coloneqq \int _a^b K(x, y) f(y) \, \mathrm{d} y. 
    \]  
    Then \(T\) is compact. More precisely, if \(\left\{ f_n \right\} \subseteq C([a, b]) \) is bounded in the supremum norm, then the sequence \(\left\{ T f_n \right\} \) has a uniformly convergent subsequence in \(C([a, b])\).    
\end{theorem}

\begin{proof}
    We divide the proof into two steps. 
    \begin{itemize}
        \item \textbf{Step 1. The operator is well-defined.} Let \(f \in C([a, b])\). By previous proposition, we khow that \(Tf \in C([a, b])\). 
        \item \textbf{Step 2. The image of every bounded sequence is relatively compact.} Let \(\left\{ f_n \right\} \subseteq C([a, b]) \) be bounded in the supremum norm. Then there exists a constant \(M > 0\) s.t. 
        \[
            \lVert f_n \rVert_\infty  \le M \quad \text{for all } n \in \mathbb{N}.  
        \]
        We shall prove the family \(\left\{ T f_n: n \in \mathbb{N} \right\} \) is uniformly bounded and equicontinuous. Then Arzela-Ascoli theorem will then imply that \(\left\{ T f_n \right\} \) has a uniformly convergent subsequence in \(C([a, b])\). 

        First, since \(K\) is continuous on the compact set \([a, b] \times [a, b]\), it is bounded. Hence, there exists \(C > 0\) s.t. 
        \[
            \vert K(x, y) \vert \le C \quad \text{for all } (x, y) \in [a, b] \times [a, b].  
        \]     
        Therefore, for every \(x \in [a, b]\), 
        \begin{align*}
            \vert (Tf_n)(x) \vert &= \left\vert \int _a^b K(x, y) f_n(y) \, \mathrm{d} y  \right\vert \le \int _a^b \vert K(x, y) \vert \vert f_n(y) \vert \, \mathrm{d} y \le \int _a^b CM \, \mathrm{d} y = CM(b - a).      
        \end{align*} 
        Thus, \(\lVert Tf_n \rVert_\infty \le CM(b - a) \) for all \(n \in \mathbb{N}\). So the family \(\left\{ Tf_n \right\} \) is uniformly bounded. Next, we prove equicontinuity. Since \(K\) is uniformly continuous on \([a, b] \times [a, b]\), given \(\varepsilon > 0\), there exists \(\delta > 0\) s.t. 
        \[
            \vert x - x^{\prime}  \vert < \delta \implies \vert K(x, y) - K(x^{\prime} , y) \vert < \frac{\varepsilon}{M(b - a)} 
        \]       
        for every \(y \in [a, b]\). If \(M = 0\), then \(f_n = 0\) for every \(n\), and the claim is trivial. Hence we may assume \(M > 0\). 

        Thus, if \(\vert x - x^{\prime}  \vert < \delta \), then 
        \begin{align*}
            \vert (T f_n)(x) - (T f_n)(x^{\prime} ) \vert &= \left\vert \int _a^b \left( K(x, y) - K(x^{\prime} , y) \right) f_n(y) \, \mathrm{d} y   \right\vert \\
            &\le \int _a^b \vert K(x, y) - K(x^{\prime} , y) \vert \vert f_n(y) \vert \, \mathrm{d} y \\
            &< \int _a^b \frac{\varepsilon}{M(b - a)} M \, \mathrm{d} y = \varepsilon.     
        \end{align*}      
        The estimate is independent of \(n\). Hence the family \(\left\{ T f_n \right\} \) is equicontinuous, and we're done.  
    \end{itemize}
\end{proof}

\chapter{The Stone-Weierstrass Theorem}
In this chapter, we discuss an general approximation theorem that has many applications to approximation problems. The theorem gives a simple and useful criterion for deciding when all continuous functions on a compact metric space can be uniformly approximated by elements of a given algebra of functions.

As consequences, we shall recover approximation results for polynomials in several variables and for trigonometric polynomials.

\section{Algebras and vector lattices}
\begin{definition}[Algebra of continuous functions]
    Let \(X\) be a compact metric space. A subset \(A \subseteq C(X)\), where \(C(X)\) denotes the space of real-valued continuous functions on \(X\), is called an algebra if \(A\) is a vector subspace of \(C(X)\) and is closed under pointwise multiplication. That is, if \(f, g \in A\), then \(fg \in A\).        
\end{definition}

For \(f, g \in C(X)\), we define 
\[
    (f \lor g)(x) \coloneqq \max \left\{ f(x), g(x) \right\}, \quad (f \land g)(x) \coloneqq \min \left\{ f(x), g(x) \right\}. 
\] 

\begin{definition}[Vector lattice]
    A subset \(L \subseteq C(X)\) is called a vector lattice if \(L\) is a vector subspace of \(C(X)\) and is closed under the two operations
    \[
        f, g \mapsto f \lor g, \quad f, g \mapsto f \land g.
    \]
    Equivalently, wheanever \(f, g \in L\), one has 
    \[
        f \lor g \in L, \quad f \land g \in L.
    \]   
\end{definition}

The operations \(f \lor g\) and \(f \land g\) can be expressed using the absolute value:
\[
    f \lor g = \frac{1}{2} (f + g) + \frac{1}{2} \vert f - g \vert, \quad f \land g = \frac{1}{2} (f + g) - \frac{1}{2} \vert f - g \vert.
\]